\begin{table}[htbp]
\centering
\caption{Hyperparameter search space for the headline scoring model, and the value taken for each. ``Fixed'' denotes a quantity that was set before the search rather than optimised over; the candidate values shown for those rows are the alternatives present in the search code but disabled.}
\label{tab:finbert_hparams_space}
\begin{tabular}{llll}
\hline
Hyperparameter & Values considered & Selected & Stage \\
\hline
        Unfrozen encoder blocks & 1, 2, 3 & 2 & Layer sweep \\
        Learning rate & log-uniform $[10^{-5},\,10^{-3}]$ & $7.35\times10^{-5}$ & Optuna \\
        Weight decay & $[5\times10^{-4},\,5\times10^{-2}]$ & 0.094$^{\dagger}$ & Optuna \\
        Optimiser & AdamW, SGD, RMSprop & AdamW & Fixed \\
        Batch size & 8, 16 & 9 & Fixed \\
        Max sequence length & 64, 128 & 128 & Fixed \\
        Momentum & $[0.001,\,0.4]$ & 0 (unused by AdamW) & Fixed \\
        Huber $\delta$ & not searched & 1 & Fixed \\
        Epochs & not searched & 100 & Fixed \\
        Regression head & not searched & frozen \texttt{Linear(3,1)} & Fixed \\
\hline
\end{tabular}
\par\smallskip
\footnotesize $^{\dagger}$ The value carried forward, 0.094, lies outside the range the search covered; it is reported as used.
\end{table}

\begin{table}[htbp]
\centering
\caption{Departures from the published configuration made for the out-of-fold re-scoring.}
\label{tab:finbert_hparams_changes}
\begin{tabular}{llll}
\hline
Setting & Published & Used here & Reason \\
\hline
        Batch size & 9 & 64 & throughput; 170 to 490 samples/s \\
        Learning rate & $7.35\times10^{-5}$ & $1.96\times10^{-4}$ & rescaled by $\sqrt{64/9}$ \\
        Precision & fp32 & fp16 (AMP) & throughput \\
        Epoch budget & 100, fixed & $\leq$100, patience 5 & stop on each fold's own validation slice \\
        Target scaling & raw & per-fold $z$ & output sd 0.698 vs target sd 0.058 at init \\
        Output activation & $\tanh$ & linear & a standardised target cannot fit in $(-1,1)$ \\
\hline
\end{tabular}
\end{table}

\begin{table}[htbp]
\centering
\caption{Training configuration of the headline scoring model, as run.}
\label{tab:finbert_hparams}
\begin{tabular}{ll}
\hline
Setting & Value \\
\hline
        Base checkpoint & \texttt{yiyanghkust/finbert-tone} (BERT-base, 109.8M parameters) \\
        Trainable parameters & last 2 encoder blocks, 14.18M of 109.8M \\
        Regression head & frozen \texttt{Linear(3,1)} on the three tone logits \\
        Output activation & linear \\
        Target & \texttt{volume\_surge\_price}, standardised per fold on training-window statistics \\
        Loss & Huber, $\delta = 1$ \\
        Optimiser & AdamW, weight decay 0.094 \\
        Learning rate & $1.96\times10^{-4}$ \\
        Batch size & 64, mixed precision (fp16) \\
        Max sequence length & 128 tokens \\
        Epoch budget & 100 maximum, early stopping patience 5 \\
        Epoch selection & last 10\% of each fold's training window, chronological, behind the purge gap \\
        Purge gap & 1 day at every fold boundary and between train and validation \\
        Seed & 42 \\
        Score calibration & per fold, standardised by that fold's validation mean and standard deviation \\
\hline
\end{tabular}
\end{table}

\paragraph{Selection procedure.} In order:

\begin{enumerate}
\item \textbf{Depth sweep.} The number of unfrozen encoder blocks was chosen first, by
      training with 1, 2 and 3 blocks unfrozen and all other settings held constant.
      Two blocks minimised both error measures and gave the highest correlation
      (MSE 0.0028, MAE 0.0313, $r = 0.2738$), and were retained.
\item \textbf{Optuna search.} With the depth fixed, the learning rate and weight decay
      were tuned by Optuna over the ranges in Table~\ref{tab:finbert_hparams_space}.
      The optimiser, batch size and sequence length were pinned rather than searched.
      Trials were scored on a random 10\% split of a 3{,}000-headline subsample.
\item \textbf{Freeze.} The resulting configuration was carried forward unchanged. It is
      \emph{not} re-tuned per fold: one configuration is used for all six folds.
\item \textbf{Throughput adjustment.} The batch size was raised from 9 to 64 and mixed
      precision enabled, with the learning rate rescaled by $\sqrt{64/9}$ to keep the
      per-step noise scale comparable. No other value was altered for speed.
\item \textbf{Conditioning correction.} The target is standardised on each fold's
      training window and the output activation made linear. At initialisation the
      published $\tanh$ head emits values with standard deviation 0.698 against a
      target standard deviation of 0.058, with mean gradient 0.293, so the model
      begins roughly twelve target standard deviations from its optimum behind a
      largely saturated nonlinearity. On a standardised target the published
      $\delta = 1$ also becomes a genuine robust threshold, whereas on the raw target
      ($|y| < 0.25$) the linear branch of the Huber loss is unreachable and the
      objective reduces to mean squared error.
\item \textbf{Per-fold epoch selection.} The only quantity chosen inside the
      out-of-fold procedure is the number of epochs, taken as the minimum of the
      validation loss on the last 10\% of each fold's own training window, behind the
      purge gap, with patience 5.
\item \textbf{Per-fold calibration.} Each fold's scores are standardised by the mean
      and standard deviation of that fold's validation predictions, both computed on
      data preceding the scored slice.
\end{enumerate}

\paragraph{Limitation.} Steps 1 and 2 were carried out on a random split of the full
sample, so the configuration itself is not free of look-ahead even though the weights
and the epoch count are. A fully nested design would repeat the search inside each
fold's training window. Because one fixed configuration is applied unchanged to every
fold, this affects the choice of configuration rather than the fold structure, and is
reported as a limitation rather than a resolved issue.
